% =====================================================================
%  Resumen / Abstract / Keywords  (Bloque UMA obligatorio)
%  ESTADO: REDACTADO. Cifras verificadas contra output/ (2026-06-03).
%  Front matter sin numerar (report).
% =====================================================================

\chapter*{Resumen}
\addcontentsline{toc}{chapter}{Resumen}

El mantenimiento del software destina una fracción muy elevada del tiempo de
desarrollo a la comprensión del código. La complejidad cognitiva, métrica
propuesta por SonarSource, cuantifica ese esfuerzo penalizando especialmente el
anidamiento de estructuras de control. Los condicionales anidados son uno de sus
principales contribuyentes y, en muchos casos, pueden combinarse de forma segura
en una única condición conjuntiva, reduciendo así la complejidad sin alterar el
comportamiento observable.

Este trabajo formaliza las condiciones bajo las cuales esa combinación es segura
(un conjunto de cuatro precondiciones estructurales verificables sobre el árbol
de sintaxis abstracta) e implementa un prototipo determinista que las comprueba y aplica la
transformación. El método experimental se articula en torno a tres preguntas de
investigación: qué casos son combinables (RQ1), qué impacto tiene la
transformación en la complejidad cognitiva (RQ2) y si los grandes modelos de
lenguaje pueden realizarla de forma correcta y automática (RQ3).

Sobre un corpus de 126 métodos Java de proyectos de código abierto, el prototipo
identifica 111 casos elegibles y reduce la complejidad cognitiva en un total de
340 puntos (todos los casos elegibles mejoran), con una coincidencia exacta
del 100\,\% frente a SonarQube en el corpus completo. En cuanto a los modelos de
lenguaje, exhiben perfiles opuestos: uno conservador, que nunca transforma en
falso, y otro más capaz pero con algún falso positivo; ninguno iguala la
garantía ni la eficiencia del prototipo determinista. El trabajo aporta, además,
un paquete de replicación completo.

\vspace{1em}
\noindent\textbf{Palabras clave:} complejidad cognitiva, refactorización,
condicionales anidados, Java, árbol de sintaxis abstracta, JavaParser, grandes
modelos de lenguaje, Design Science Research.

\chapter*{Abstract}
\addcontentsline{toc}{chapter}{Abstract}

Software maintenance devotes a large fraction of development time to code
comprehension. Cognitive complexity, a metric proposed by SonarSource, quantifies
that effort by penalising nesting in control structures. Nested conditionals are
among its main contributors, and in many cases they can be safely combined into a
single conjunctive condition, reducing complexity without altering observable
behaviour.

This work formalises the conditions under which such combination is safe (a set
of four structural verifiable preconditions on the abstract syntax tree) and implements a
deterministic prototype that checks them and applies the transformation. The
experimental method addresses three research questions: which cases are
combinable (RQ1), what impact the transformation has on cognitive complexity
(RQ2), and whether large language models can perform it correctly and
automatically (RQ3).

On a corpus of 126 Java methods from open-source projects, the prototype
identifies 111 eligible cases and reduces cognitive complexity by a total of 340
points (all eligible cases improve), with 100\,\% exact agreement against
SonarQube on the full corpus. Regarding the language models, they exhibit
opposite profiles: one conservative, never producing a false transformation, and
one more capable but with occasional false positives; neither matches the
correctness guarantees or the efficiency of the deterministic prototype. The work
also delivers a complete replication package.

\vspace{1em}
\noindent\textbf{Keywords:} cognitive complexity, refactoring, nested
conditionals, Java, abstract syntax tree, JavaParser, large language models,
Design Science Research.
